% Esquisse d'un programme (cote 119, copy 1): archive cote 119, site page 55 (PDF page 56), typescript page 53.
% typescript. Draft compared: EsquisseFr.pdf (text layer):1888-1936.
% AI-transcribed by Claude Opus 5.5 (claude-opus-5-5, adaptive thinking, effort high), Message Batches API, batch esq1 (msgbatch_01LNsogFyF3y6KSp9UbtJtBk), 2026-11-12; not yet checked by eye.
\scanpage{119/55}
\mspage{53}

la définition du cercle. Celle-ci m'avait impréssionné par sa sim-%
plicité et son évidence, alors que la propriété de ``rotondité
parfaite'' du cercle m'apparaissait auparavant comme une réalité mysté-%
rieuse au delà des mots. C'est à ce moment, je crois, que j'ai
entrevu pour la première fois (sans bien sûr me le formuler en
ces termes) la puissance créatrice d'une ``bonne'' définition ma-%
thématique, d'une \uline{formulation} qui décrit l'essence. Aujourd'hui
encore, il semble que la fascination qu'a exercé sur moi cette
puissance-là n'a rien perdu de sa force.

(\textsuperscript{3}) Plus généralement, au delà des variétés dites ``anabéliennes''
sur des corps de type fini, la géométrie algébrique anabélienne
(telle qu'elle s'est dégagé il y a quelques années) amène à une
description, en termes de groupes profinis uniquement, de la ca-%
tégorie des schémas de type fini sur la base absolue $\mathbb{Q}$ (voir
même $\mathbb{Z}$), et par là même, en principe, de la catégorie des sché-%
mas quelconques (par des passages à la limite convenables). Il
s'agit donc d'une construction ``qui fait semblant'' d'ignorer les
anneaux (tels que $\mathbb{Q}$, les algèbres de type fini sur $\mathbb{Q}$ etc) et
les équations algébriques qui servent traditionnellement à décri-%
re les schémas, en travaillant directement avec leurs topos éta-%
les, exprimables en termes de systèmes de groupes profinis. Un
grain de sel cependant : pour pouvoir espérer reconstituer un
schéma (de type fini sur $\mathbb{Q}$ disons) à partir de son topos étale,
qui est un invariant purement topologique, il convient de se pla-%
cer, non dans la catégorie des schémas (de type fini sur $\mathbb{Q}$ en
l'occurence), mais dans celle qui s'en déduit par ``localisation'',
en rendant inversibles les morphismes qui sont des ``homéomorphismes
universels'', i.e. qui sont finis, radiciels et surjectifs. Le
développement d'une telle traduction d'un ``monde géométrique''
(savoir celui des schémas, multiplicités schématiques etc) en
termes d'un ``monde algébrique'' (celui des groupes profinis, et
systèmes de groupes profinis décrivant des topos (dits ``étales'')
convenables) peut être considéré comme un aboutissement ultime

\margin{55}
